\answer
\begin{enumerate}

%% 1章 問1（1.1節）
\item[1.1節 問1]
1秒間に通過する電荷は$Q = It = 10\,\mathrm{C}$。
電子1個の電荷の大きさは$e = 1.602\times10^{-19}$\,Cだから，
\begin{equation*}
N = \frac{Q}{e} = \frac{10}{1.602\times10^{-19}} \approx 6.2\times10^{19}\,\text{個}
\end{equation*}
毎秒6千京個もの電子が通過している。電流はキャリアの莫大な数の流れである。

%% 1章 問2（1.2節）
\item[1.2節 問2]
リン（電子15個）：$1s^2\,2s^2\,2p^6\,3s^2\,3p^3$で価電子は5個。
ホウ素（電子5個）：$1s^2\,2s^2\,2p^1$で価電子は3個。
4価のシリコンに対して価電子が1個多いリンはドナーに，
1個少ないホウ素はアクセプタになる（\ref{sec1.4}節）。

%% 1章 問3（1.3節）
\item[1.3節 問3]
$\dfrac{E-E_F}{kT} = \dfrac{0.55}{0.026} \approx 21$より
$\exp(21) \gg 1$だから，分母の1は無視できて
\begin{equation*}
f(E) \approx \exp\left(-\frac{E-E_F}{kT}\right) = \exp(-21) \approx 7\times10^{-10}
\end{equation*}
10億分の1以下というきわめて小さい確率だが，ゼロではない。
価電子帯には$10^{22}\,\mathrm{cm^{-3}}$規模の電子があるから，
これだけ小さい確率でも室温の真性キャリア濃度（$10^{10}\,\mathrm{cm^{-3}}$規模）が生まれる。

%% 1章 問4（1.4節）
\item[1.4節 問4]
アクセプタをドーピングしたのでp形半導体である。
多数キャリアは正孔で$p \approx N_a = 10^{17}\,\mathrm{cm^{-3}}$。
少数キャリアは電子で，質量作用の法則\ref{eq1.3}より
\begin{equation*}
n = \frac{n_i^2}{p} = \frac{(1.5\times10^{10})^2}{10^{17}}
\approx 2.3\times10^{3}\,\mathrm{cm^{-3}}
\end{equation*}

%% 1章 問5（1.5節）
\item[1.5節 問5]
式\ref{eq1.4}の右辺：
$\mu\mathcal{E}$の単位は
$\dfrac{\mathrm{m^2}}{\mathrm{V\cdot s}}\times\dfrac{\mathrm{V}}{\mathrm{m}}
= \mathrm{m/s}$となり，左辺（速度）と一致する。
式\ref{eq1.5}の右辺：
$en\mu\mathcal{E}$の単位は
$\mathrm{C}\times\mathrm{m^{-3}}\times\dfrac{\mathrm{m^2}}{\mathrm{V\cdot s}}\times\dfrac{\mathrm{V}}{\mathrm{m}}
= \dfrac{\mathrm{C}}{\mathrm{s\cdot m^2}} = \mathrm{A/m^2}$となり，
左辺（電流密度）と一致する。

%% 1章 問6（1.6節）
\item[1.6節 問6]
根号の中の単位は
\begin{equation*}
\frac{\mathrm{F/m}\times\mathrm{V}}{\mathrm{C}}\times\mathrm{m^3}
= \frac{\mathrm{F\cdot V}}{\mathrm{C}}\times\mathrm{m^2}
\end{equation*}
$1\,\mathrm{F} = 1$\,C/Vより$\mathrm{F\cdot V/C} = 1$（無次元）だから，
根号の中は$\mathrm{m^2}$，したがって右辺の単位はmとなり，
左辺（空乏層幅）と一致する。

%% 1章 問7（1.7節）
\item[1.7節 問7]
n形半導体で$\phi_m = 5.1\,\mathrm{eV} > \phi_s = 4.3$\,eVだから，
整流性のショットキー接合になる（\ref{tab1.1}）。
ショットキー障壁の高さは式\ref{eq1.12}より
\begin{equation*}
\phi_B = \phi_m - \chi_s = 5.1 - 4.05 \approx 1.1\,\mathrm{eV}
\end{equation*}

\end{enumerate}
